\section{Related Work}
\label{sec:related}

\paragraph{Automated code and harness optimization.}

A long line of work, encompassing OPRO, FunSearch, AlphaEvolve, and ShinkaEvolve, has applied LLMs within larger search scaffolds to discover novel or optimized programs 
\citep{yang2024opro,romera2024funsearch,novikov2025alphaevolve,lange2025shinkaevolveopenendedsampleefficientprogram}. Existing approaches to \emph{harness optimization} differ along two axes: which parts of the harness they may
modify, and the role assigned to the LLM during search. Prompt optimization
methods such as DSPy \citep{khattab2024dspy} optimize prompts while holding the surrounding program fixed. Meta Context Engineering
\citep{ye2026metacontextengineeringagentic} searches over \emph{skills} and context
artifacts in a bi-level procedure. TextGrad
\citep{yuksekgonul2025textgrad}, Trace  \citep{cheng2024trace}, and LLM-AutoDiff
\citep{yin2025adalflow} optimize arbitrary text components in chained workflows by
back-propagating textual feedback using LLMs. GEPA \citep{agrawal2026gepa} optimizes textual artifacts by using
reflective feedback to guide an evolutionary process.  A second line expands the search space to the
agent program itself: STOP \citep{zelikman2024stop} recursively improves a
scaffolding program. ADAS \citep{hu2025adas} searches over agent designs
represented as code. AFlow \citep{zhang2025aflow} performs tree search over
code-defined workflows. Darwin
G\"odel Machine \citep{zhang2026dgm}, G\"odel Agent
\citep{yin2025godelagent}, and SICA
\citep{robeyns2025sica} study recursive \emph{self-modification}.

Closest to the regime \heb\ measures are optimizers that act end-to-end:
\citet{ursekar2026vero} and \citet{lee2026metaharnessendtoendoptimizationmodel} let a coding agent read the target source, inspect
scores and execution traces from prior candidates, and choose what evidence to
gather, rather than orchestrating mutations in a fixed
search algorithm. Other
systems in this regime include Agentic Harness Engineering
\citep{lin2026agenticharness} and HarnessX
\citep{chen2026harnessx}. Several of these works contribute optimization
\emph{methods}, evaluated on method-specific seeds, budgets, search
spaces, and scoring protocols, which makes their reported outcomes difficult to
compare. \heb\ is complementary: it fixes the optimization problem and the
evaluation protocol so that optimizer models, harnesses, and search algorithms can
be compared on common ground. \heb\ builds upon the infrastructure and protocol in 
\vero\ and contributes a suite of optimization tasks. 

\paragraph{Coding agent benchmarks.}
Coding agent evaluation has grown from function-level synthesis
\citep{chen2021humaneval,austin2021mbpp} to repository-scale tasks in real
execution environments \citep{jimenez2024swebench}, and tasks requiring iterative optimization, such as machine learning
engineering \citep{chan2025mlebench} and kernel optimization \citep{ouyang2025kernelbench}. \heb\ similarly requires coding agents to navigate the target system codebase and iterate on environmental feedback; it differs in what the optimization target is (i.e. a harness), the noisiness of evaluation feedback, and the types of bounds imposed on its search. 


\begin{figure*}[t]
\centering
\includegraphics[width=\linewidth]{figures/R01_lm_showdown.pdf}
\caption{\textbf{Optimizer models separate more than their coding harnesses.}
\emph{Left:} normalized gain for every run in the controlled two-harness design.
Marker shape denotes the harness; horizontal offsets prevent overlap and carry
no task meaning. \emph{Right:} LSS-$\lambda$, estimated from
the balanced shared-harness grid on the three competent-seed tasks.}
\label{fig:showdown}
\end{figure*}
